\subsection{Przypadki funkcjonalne}
\subsubsection{TC-F01 – Rejestracja konta poprawnymi danymi}
{\small
\begin{longtable}{p{\dimexpr 0.270\linewidth-2\tabcolsep\relax}p{\dimexpr 0.730\linewidth-2\tabcolsep\relax}}
\toprule
\textbf{Pole} & \textbf{Wartość} \\
\midrule
\endhead
Priorytet & Wysoki \\
Technika & Podział na klasy równoważności (klasa poprawna: unikalna nazwa, poprawny e-mail, hasło $\geq$ 8 znaków) \\
Typ testu & Funkcjonalny \\
Poziom testu & Testowanie systemowe (UI, Selenium) \\
Warunki wstępne / dane & Rejestracja włączona; nazwa \texttt{reg-<losowy>}, e-mail \texttt{<nazwa>@example.com}, hasło \texttt{Valid123!pass}; użytkownik niezalogowany \\
Powiązane funkcje & F01 \\
\bottomrule
\end{longtable}
}
{\small
\begin{longtable}{p{\dimexpr 0.050\linewidth-2\tabcolsep\relax}p{\dimexpr 0.480\linewidth-2\tabcolsep\relax}p{\dimexpr 0.470\linewidth-2\tabcolsep\relax}}
\toprule
\textbf{\#} & \textbf{Krok} & \textbf{Oczekiwany rezultat} \\
\midrule
\endhead
1 & Otwórz \texttt{/user/sign\_up} & Widoczny formularz z polami: nazwa, e-mail, hasło, powtórz hasło \\
2 & Wypełnij pola danymi testowymi & Pola przyjmują wartości, brak komunikatów walidacji \\
3 & Kliknij „Register Account” & Przekierowanie na stronę główną; w nawigacji awatar z nazwą nowego użytkownika \\
4 & Sprawdź \texttt{GET /api/v1/users/<nazwa>} & HTTP 200 \\
\bottomrule
\end{longtable}
}
Warunki końcowe: konto istnieje i jest aktywne, użytkownik zalogowany; po teście konto usunięte przez administratora.

\subsubsection{TC-F02 – Walidacja minimalnej długości hasła przy rejestracji}
{\small
\begin{longtable}{p{\dimexpr 0.270\linewidth-2\tabcolsep\relax}p{\dimexpr 0.730\linewidth-2\tabcolsep\relax}}
\toprule
\textbf{Pole} & \textbf{Wartość} \\
\midrule
\endhead
Priorytet & Wysoki \\
Technika & Analiza wartości brzegowych (dwupunktowa: 7 i 8 znaków; granica \texttt{MIN\_PASSWORD\_LENGTH = 8}) \\
Typ testu & Funkcjonalny \\
Poziom testu & Testowanie systemowe (UI, Selenium) \\
Warunki wstępne / dane & Użytkownik niezalogowany; hasło \texttt{Ab1!Ab1} (7 zn.) oraz \texttt{Ab1!Ab1!} (8 zn.) \\
Powiązane funkcje & F01 \\
\bottomrule
\end{longtable}
}
{\small
\begin{longtable}{p{\dimexpr 0.050\linewidth-2\tabcolsep\relax}p{\dimexpr 0.480\linewidth-2\tabcolsep\relax}p{\dimexpr 0.470\linewidth-2\tabcolsep\relax}}
\toprule
\textbf{\#} & \textbf{Krok} & \textbf{Oczekiwany rezultat} \\
\midrule
\endhead
1 & Otwórz \texttt{/user/sign\_up}, wypełnij formularz hasłem 7-znakowym, wyślij & Komunikat „Password length cannot be less than 8 characters.”; brak przekierowania \\
2 & Sprawdź \texttt{GET /api/v1/users/<nazwa>} & HTTP 404 – konto nie powstało \\
3 & Powtórz krok 1 z hasłem 8-znakowym (nowa nazwa) & Konto utworzone, użytkownik zalogowany (awatar w nawigacji) \\
\bottomrule
\end{longtable}
}
Warunki końcowe: istnieje wyłącznie konto z hasłem 8-znakowym; usunięte po teście.

\subsubsection{TC-F03 – Logowanie poprawnymi danymi}
{\small
\begin{longtable}{p{\dimexpr 0.270\linewidth-2\tabcolsep\relax}p{\dimexpr 0.730\linewidth-2\tabcolsep\relax}}
\toprule
\textbf{Pole} & \textbf{Wartość} \\
\midrule
\endhead
Priorytet & Wysoki \\
Technika & Podział na klasy równoważności (klasa poprawna) \\
Typ testu & Funkcjonalny \\
Poziom testu & Testowanie systemowe (UI, Selenium) \\
Warunki wstępne / dane & Istnieje aktywny użytkownik testowy; przeglądarka bez sesji \\
Powiązane funkcje & F02 \\
\bottomrule
\end{longtable}
}
{\small
\begin{longtable}{p{\dimexpr 0.050\linewidth-2\tabcolsep\relax}p{\dimexpr 0.480\linewidth-2\tabcolsep\relax}p{\dimexpr 0.470\linewidth-2\tabcolsep\relax}}
\toprule
\textbf{\#} & \textbf{Krok} & \textbf{Oczekiwany rezultat} \\
\midrule
\endhead
1 & Otwórz \texttt{/user/login} & Formularz z polami nazwy i hasła \\
2 & Wpisz poprawną nazwę i hasło, kliknij „Sign In” & Przekierowanie na \texttt{/} \\
3 & Sprawdź nawigację & Awatar z tytułem równym nazwie użytkownika \\
\bottomrule
\end{longtable}
}
Warunki końcowe: aktywna sesja użytkownika testowego.

\subsubsection{TC-F04 – Odrzucenie logowania z błędnym hasłem}
{\small
\begin{longtable}{p{\dimexpr 0.270\linewidth-2\tabcolsep\relax}p{\dimexpr 0.730\linewidth-2\tabcolsep\relax}}
\toprule
\textbf{Pole} & \textbf{Wartość} \\
\midrule
\endhead
Priorytet & Wysoki \\
Technika & Podział na klasy równoważności (klasa niepoprawna: błędne hasło) \\
Typ testu & Funkcjonalny \\
Poziom testu & Testowanie systemowe (UI, Selenium) \\
Warunki wstępne / dane & Istnieje użytkownik testowy; hasło \texttt{wrong-password} \\
Powiązane funkcje & F02 \\
\bottomrule
\end{longtable}
}
{\small
\begin{longtable}{p{\dimexpr 0.050\linewidth-2\tabcolsep\relax}p{\dimexpr 0.480\linewidth-2\tabcolsep\relax}p{\dimexpr 0.470\linewidth-2\tabcolsep\relax}}
\toprule
\textbf{\#} & \textbf{Krok} & \textbf{Oczekiwany rezultat} \\
\midrule
\endhead
1 & Otwórz \texttt{/user/login} & Formularz logowania \\
2 & Wpisz poprawną nazwę i błędne hasło, wyślij & Komunikat „Username or password is incorrect.” \\
3 & Sprawdź URL i nawigację & URL zawiera \texttt{/user/login}; brak awatara (użytkownik niezalogowany) \\
\bottomrule
\end{longtable}
}
Warunki końcowe: brak sesji; konto nie zostało zablokowane.

\subsubsection{TC-F05 – Utworzenie repozytorium z plikiem README}
{\small
\begin{longtable}{p{\dimexpr 0.270\linewidth-2\tabcolsep\relax}p{\dimexpr 0.730\linewidth-2\tabcolsep\relax}}
\toprule
\textbf{Pole} & \textbf{Wartość} \\
\midrule
\endhead
Priorytet & Wysoki \\
Technika & Podział na klasy równoważności (poprawna nazwa, repo publiczne, inicjalizacja włączona) \\
Typ testu & Funkcjonalny \\
Poziom testu & Testowanie systemowe (UI, Selenium + weryfikacja API) \\
Warunki wstępne / dane & Zalogowany użytkownik testowy; nazwa \texttt{ui-repo-<losowy>}, opis „Repo z testu Selenium” \\
Powiązane funkcje & F03, F04 \\
\bottomrule
\end{longtable}
}
{\small
\begin{longtable}{p{\dimexpr 0.050\linewidth-2\tabcolsep\relax}p{\dimexpr 0.480\linewidth-2\tabcolsep\relax}p{\dimexpr 0.470\linewidth-2\tabcolsep\relax}}
\toprule
\textbf{\#} & \textbf{Krok} & \textbf{Oczekiwany rezultat} \\
\midrule
\endhead
1 & Otwórz \texttt{/repo/create} & Formularz nowego repozytorium \\
2 & Podaj nazwę i opis, zaznacz „Initialize Repository”, kliknij „Create Repository” & Przekierowanie na \texttt{/<użytkownik>/<nazwa>} \\
3 & Sprawdź listę plików & Lista zawiera \texttt{README.md} \\
4 & Sprawdź sekcję README & Zawiera nazwę repo i opis \\
5 & \texttt{GET /api/v1/repos/<użytkownik>/<nazwa>} & \texttt{description} = opis, \texttt{private} = false \\
\bottomrule
\end{longtable}
}
Warunki końcowe: repozytorium istnieje z jednym commitem; usunięte po teście.

\subsubsection{TC-F06 – Odrzucenie duplikatu nazwy repozytorium}
{\small
\begin{longtable}{p{\dimexpr 0.270\linewidth-2\tabcolsep\relax}p{\dimexpr 0.730\linewidth-2\tabcolsep\relax}}
\toprule
\textbf{Pole} & \textbf{Wartość} \\
\midrule
\endhead
Priorytet & Średni \\
Technika & Zgadywanie błędów (error guessing) \\
Typ testu & Funkcjonalny \\
Poziom testu & Testowanie systemowe (UI, Selenium) \\
Warunki wstępne / dane & Zalogowany użytkownik; istnieje jego repozytorium o nazwie X (utworzone przez API) \\
Powiązane funkcje & F03 \\
\bottomrule
\end{longtable}
}
{\small
\begin{longtable}{p{\dimexpr 0.050\linewidth-2\tabcolsep\relax}p{\dimexpr 0.480\linewidth-2\tabcolsep\relax}p{\dimexpr 0.470\linewidth-2\tabcolsep\relax}}
\toprule
\textbf{\#} & \textbf{Krok} & \textbf{Oczekiwany rezultat} \\
\midrule
\endhead
1 & Otwórz \texttt{/repo/create}, podaj nazwę X, wyślij & Komunikat „The repository name is already used.” \\
2 & Sprawdź URL & Pozostaje \texttt{/repo/create} \\
3 & \texttt{GET /api/v1/user/repos} & Nazwa X występuje dokładnie raz \\
\bottomrule
\end{longtable}
}
Warunki końcowe: stan repozytoriów użytkownika bez zmian.

\subsubsection{TC-F07 – Utworzenie zgłoszenia z etykietą}
{\small
\begin{longtable}{p{\dimexpr 0.270\linewidth-2\tabcolsep\relax}p{\dimexpr 0.730\linewidth-2\tabcolsep\relax}}
\toprule
\textbf{Pole} & \textbf{Wartość} \\
\midrule
\endhead
Priorytet & Wysoki \\
Technika & Testowanie oparte na kryteriach akceptacji \\
Typ testu & Funkcjonalny \\
Poziom testu & Testowanie systemowe (UI, Selenium + weryfikacja API) \\
Warunki wstępne / dane & Zalogowany właściciel repozytorium; w repo istnieje etykieta \texttt{bug}; tytuł „Błąd logowania-<losowy>” \\
Powiązane funkcje & F05, F09 \\
\bottomrule
\end{longtable}
}
{\small
\begin{longtable}{p{\dimexpr 0.050\linewidth-2\tabcolsep\relax}p{\dimexpr 0.480\linewidth-2\tabcolsep\relax}p{\dimexpr 0.470\linewidth-2\tabcolsep\relax}}
\toprule
\textbf{\#} & \textbf{Krok} & \textbf{Oczekiwany rezultat} \\
\midrule
\endhead
1 & Otwórz \texttt{/<repo>/issues/new} & Formularz z polem tytułu, opisem i panelem bocznym \\
2 & Wpisz tytuł i opis & Wartości widoczne w polach \\
3 & W panelu „Labels” wybierz \texttt{bug}, zamknij listę & Etykieta \texttt{bug} widoczna w panelu \\
4 & Kliknij „Create Issue” & Przekierowanie na \texttt{/<repo>/issues/<n>}; tytuł i etykieta \texttt{bug} widoczne \\
5 & \texttt{GET /api/v1/repos/<repo>/issues/<n>} & \texttt{state} = open, \texttt{labels} = [\texttt{bug}] \\
\bottomrule
\end{longtable}
}
Warunki końcowe: w repo istnieje 1 otwarte zgłoszenie z etykietą \texttt{bug}.

\subsubsection{TC-F08 – Komentarz, zamknięcie i ponowne otwarcie zgłoszenia}
{\small
\begin{longtable}{p{\dimexpr 0.270\linewidth-2\tabcolsep\relax}p{\dimexpr 0.730\linewidth-2\tabcolsep\relax}}
\toprule
\textbf{Pole} & \textbf{Wartość} \\
\midrule
\endhead
Priorytet & Wysoki \\
Technika & Testowanie przejść między stanami (open $\rightarrow$ closed $\rightarrow$ open; przejścia „zamknij z komentarzem”, „otwórz ponownie”) \\
Typ testu & Funkcjonalny \\
Poziom testu & Testowanie systemowe (UI, Selenium + weryfikacja API) \\
Warunki wstępne / dane & Zalogowany właściciel; otwarte zgłoszenie „Do zamknięcia” utworzone przez API \\
Powiązane funkcje & F05 \\
\bottomrule
\end{longtable}
}
{\small
\begin{longtable}{p{\dimexpr 0.050\linewidth-2\tabcolsep\relax}p{\dimexpr 0.480\linewidth-2\tabcolsep\relax}p{\dimexpr 0.470\linewidth-2\tabcolsep\relax}}
\toprule
\textbf{\#} & \textbf{Krok} & \textbf{Oczekiwany rezultat} \\
\midrule
\endhead
1 & Otwórz stronę zgłoszenia & Przycisk statusu „Close Issue” \\
2 & Dodaj komentarz „Potwierdzam błąd” & Komentarz widoczny na osi czasu; stan nadal open \\
3 & Wpisz „Naprawione” i kliknij „Close with Comment” & Stan w API = closed; przycisk zmienia się na „Reopen Issue” \\
4 & Kliknij „Reopen Issue” & Stan w API = open \\
5 & \texttt{GET .../issues/<n>/comments} & Dokładnie 2 komentarze w kolejności: „Potwierdzam błąd”, „Naprawione” \\
\bottomrule
\end{longtable}
}
Warunki końcowe: zgłoszenie otwarte, 2 komentarze, 2 zdarzenia zmiany stanu na osi czasu.

\subsubsection{TC-F09 – Edycja pliku w edytorze przeglądarkowym i commit}
{\small
\begin{longtable}{p{\dimexpr 0.270\linewidth-2\tabcolsep\relax}p{\dimexpr 0.730\linewidth-2\tabcolsep\relax}}
\toprule
\textbf{Pole} & \textbf{Wartość} \\
\midrule
\endhead
Priorytet & Średni \\
Technika & Testowanie oparte na kryteriach akceptacji \\
Typ testu & Funkcjonalny \\
Poziom testu & Testowanie systemowe (UI, Selenium + weryfikacja API) \\
Warunki wstępne / dane & Zalogowany właściciel repo z \texttt{README.md}; nowa treść „Treść zmieniona w edytorze”; opis commita „Aktualizacja README z UI” \\
Powiązane funkcje & F08, F04 \\
\bottomrule
\end{longtable}
}
{\small
\begin{longtable}{p{\dimexpr 0.050\linewidth-2\tabcolsep\relax}p{\dimexpr 0.480\linewidth-2\tabcolsep\relax}p{\dimexpr 0.470\linewidth-2\tabcolsep\relax}}
\toprule
\textbf{\#} & \textbf{Krok} & \textbf{Oczekiwany rezultat} \\
\midrule
\endhead
1 & Otwórz \texttt{/<repo>/\_edit/main/README.md} & Edytor kodu z bieżącą treścią pliku \\
2 & Zaznacz całą treść i zastąp nową & Edytor zawiera wyłącznie nową treść \\
3 & Wpisz opis commita, wybierz „Commit directly”, kliknij „Commit Changes” & Przekierowanie do widoku pliku \texttt{/src/branch/main/...} \\
4 & \texttt{GET .../contents/README.md} & Zdekodowana treść = nowa treść \\
5 & \texttt{GET .../commits?limit=1} & Wiadomość ostatniego commita = podany opis \\
\bottomrule
\end{longtable}
}
Warunki końcowe: gałąź \texttt{main} ma nowy commit ze zmienionym README.

\subsubsection{TC-F10 – Utworzenie i scalenie pull requesta}
{\small
\begin{longtable}{p{\dimexpr 0.270\linewidth-2\tabcolsep\relax}p{\dimexpr 0.730\linewidth-2\tabcolsep\relax}}
\toprule
\textbf{Pole} & \textbf{Wartość} \\
\midrule
\endhead
Priorytet & Wysoki \\
Technika & Testowanie oparte na kryteriach akceptacji (przepływ end-to-end) \\
Typ testu & Funkcjonalny \\
Poziom testu & Testowanie systemowe (UI, Selenium + weryfikacja API) \\
Warunki wstępne / dane & Zalogowany właściciel; gałąź \texttt{feature-x} z plikiem \texttt{feature.txt} (utworzona przez API); tytuł PR „Dodaje feature.txt” \\
Powiązane funkcje & F06 \\
\bottomrule
\end{longtable}
}
{\small
\begin{longtable}{p{\dimexpr 0.050\linewidth-2\tabcolsep\relax}p{\dimexpr 0.480\linewidth-2\tabcolsep\relax}p{\dimexpr 0.470\linewidth-2\tabcolsep\relax}}
\toprule
\textbf{\#} & \textbf{Krok} & \textbf{Oczekiwany rezultat} \\
\midrule
\endhead
1 & Otwórz \texttt{/<repo>/compare/main...feature-x} & Diff z 1 zmienionym plikiem, przycisk „New Pull Request” \\
2 & Kliknij „New Pull Request”, wpisz tytuł, „Create Pull Request” & Przekierowanie na \texttt{/<repo>/pulls/<n>} \\
3 & Kliknij „Create merge commit” i potwierdź & Etykieta stanu „Merged” \\
4 & \texttt{GET .../pulls/<n>} & \texttt{merged} = true \\
5 & \texttt{GET .../contents/feature.txt?ref=main} & HTTP 200 – plik w \texttt{main} \\
\bottomrule
\end{longtable}
}
Warunki końcowe: PR scalony i zamknięty, \texttt{main} zawiera \texttt{feature.txt}.

\subsubsection{TC-F11 – Wyszukiwanie repozytorium po nazwie}
{\small
\begin{longtable}{p{\dimexpr 0.270\linewidth-2\tabcolsep\relax}p{\dimexpr 0.730\linewidth-2\tabcolsep\relax}}
\toprule
\textbf{Pole} & \textbf{Wartość} \\
\midrule
\endhead
Priorytet & Średni \\
Technika & Podział na klasy równoważności (repo pasujące / niepasujące do frazy) \\
Typ testu & Funkcjonalny \\
Poziom testu & Testowanie systemowe (UI, Selenium) \\
Warunki wstępne / dane & Zalogowany użytkownik; repo \texttt{szukaj-<losowy>-app} oraz repo o niepasującej nazwie \\
Powiązane funkcje & F10 \\
\bottomrule
\end{longtable}
}
{\small
\begin{longtable}{p{\dimexpr 0.050\linewidth-2\tabcolsep\relax}p{\dimexpr 0.480\linewidth-2\tabcolsep\relax}p{\dimexpr 0.470\linewidth-2\tabcolsep\relax}}
\toprule
\textbf{\#} & \textbf{Krok} & \textbf{Oczekiwany rezultat} \\
\midrule
\endhead
1 & Otwórz \texttt{/explore/repos} & Pole wyszukiwania i lista repozytoriów \\
2 & Wpisz frazę \texttt{szukaj-<losowy>}, Enter & URL zawiera \texttt{q=} \\
3 & Odczytaj wyniki & Repo pasujące jest na liście, repo niepasujące – nie \\
\bottomrule
\end{longtable}
}
Warunki końcowe: brak zmian danych (oba repo usunięte po teście).

\subsubsection{TC-F12 – Wylogowanie blokuje dostęp do ustawień konta}
{\small
\begin{longtable}{p{\dimexpr 0.270\linewidth-2\tabcolsep\relax}p{\dimexpr 0.730\linewidth-2\tabcolsep\relax}}
\toprule
\textbf{Pole} & \textbf{Wartość} \\
\midrule
\endhead
Priorytet & Wysoki \\
Technika & Testowanie przejść między stanami (zalogowany $\rightarrow$ niezalogowany) \\
Typ testu & Funkcjonalny \\
Poziom testu & Testowanie systemowe (UI, Selenium) \\
Warunki wstępne / dane & Zalogowany użytkownik testowy \\
Powiązane funkcje & F02, F07 \\
\bottomrule
\end{longtable}
}
{\small
\begin{longtable}{p{\dimexpr 0.050\linewidth-2\tabcolsep\relax}p{\dimexpr 0.480\linewidth-2\tabcolsep\relax}p{\dimexpr 0.470\linewidth-2\tabcolsep\relax}}
\toprule
\textbf{\#} & \textbf{Krok} & \textbf{Oczekiwany rezultat} \\
\midrule
\endhead
1 & Otwórz menu awatara, kliknij „Sign Out” & W nawigacji link „Sign In”, brak awatara \\
2 & Otwórz \texttt{/user/settings} & Przekierowanie na \texttt{/user/login?redirect\_to=\%2Fuser\%2Fsettings} \\
\bottomrule
\end{longtable}
}
Warunki końcowe: brak aktywnej sesji.

\subsubsection{TC-F13 – CRUD repozytorium przez REST API}
{\small
\begin{longtable}{p{\dimexpr 0.270\linewidth-2\tabcolsep\relax}p{\dimexpr 0.730\linewidth-2\tabcolsep\relax}}
\toprule
\textbf{Pole} & \textbf{Wartość} \\
\midrule
\endhead
Priorytet & Wysoki \\
Technika & Testowanie oparte na kryteriach akceptacji (cykl życia zasobu) \\
Typ testu & Funkcjonalny \\
Poziom testu & Testowanie systemowe (REST API) \\
Warunki wstępne / dane & Użytkownik testowy z uwierzytelnieniem Basic; nazwa \texttt{crud-<losowy>} \\
Powiązane funkcje & F03 \\
\bottomrule
\end{longtable}
}
{\small
\begin{longtable}{p{\dimexpr 0.050\linewidth-2\tabcolsep\relax}p{\dimexpr 0.480\linewidth-2\tabcolsep\relax}p{\dimexpr 0.470\linewidth-2\tabcolsep\relax}}
\toprule
\textbf{\#} & \textbf{Krok} & \textbf{Oczekiwany rezultat} \\
\midrule
\endhead
1 & \texttt{POST /user/repos} z opisem „v1” & 201; \texttt{full\_name} = \texttt{<user>/<nazwa>} \\
2 & \texttt{GET /repos/<user>/<nazwa>} & 200; \texttt{description} = „v1” \\
3 & \texttt{PATCH} z \texttt{description} „v2”, \texttt{has\_wiki} false & 200; pola zaktualizowane \\
4 & \texttt{DELETE /repos/<user>/<nazwa>} & 204 \\
5 & \texttt{GET /repos/<user>/<nazwa>} & 404 \\
\bottomrule
\end{longtable}
}
Warunki końcowe: repozytorium nie istnieje.

\subsubsection{TC-F14 – Widoczność repozytorium wg roli (tablica decyzyjna)}
{\small
\begin{longtable}{p{\dimexpr 0.270\linewidth-2\tabcolsep\relax}p{\dimexpr 0.730\linewidth-2\tabcolsep\relax}}
\toprule
\textbf{Pole} & \textbf{Wartość} \\
\midrule
\endhead
Priorytet & Wysoki \\
Technika & Testowanie z użyciem tablicy decyzyjnej (warunki: widoczność $\times$ rola; akcja: kod HTTP) \\
Typ testu & Funkcjonalny (bezpieczeństwo – kontrola dostępu) \\
Poziom testu & Testowanie systemowe (REST API) \\
Warunki wstępne / dane & Właściciel, drugi użytkownik, klient anonimowy; repo publiczne lub prywatne utworzone dla każdej reguły \\
Powiązane funkcje & F07 \\
\bottomrule
\end{longtable}
}
Tablica decyzyjna (6 reguł, każda to osobne wykonanie):

{\small
\begin{longtable}{p{\dimexpr 0.143\linewidth-2\tabcolsep\relax}p{\dimexpr 0.143\linewidth-2\tabcolsep\relax}p{\dimexpr 0.143\linewidth-2\tabcolsep\relax}p{\dimexpr 0.143\linewidth-2\tabcolsep\relax}p{\dimexpr 0.143\linewidth-2\tabcolsep\relax}p{\dimexpr 0.143\linewidth-2\tabcolsep\relax}p{\dimexpr 0.143\linewidth-2\tabcolsep\relax}}
\toprule
\textbf{Reguła} & \textbf{R1} & \textbf{R2} & \textbf{R3} & \textbf{R4} & \textbf{R5} & \textbf{R6} \\
\midrule
\endhead
Repo prywatne & N & N & N & T & T & T \\
Rola & właściciel & inny & anonim & właściciel & inny & anonim \\
**Oczekiwany kod \texttt{GET /repos/...}** & 200 & 200 & 200 & 200 & \textbf{404} & \textbf{404} \\
\bottomrule
\end{longtable}
}
{\small
\begin{longtable}{p{\dimexpr 0.050\linewidth-2\tabcolsep\relax}p{\dimexpr 0.480\linewidth-2\tabcolsep\relax}p{\dimexpr 0.470\linewidth-2\tabcolsep\relax}}
\toprule
\textbf{\#} & \textbf{Krok} & \textbf{Oczekiwany rezultat} \\
\midrule
\endhead
1 & Utwórz repo o widoczności z reguły & 201 \\
2 & Wykonaj \texttt{GET /repos/<repo>} jako rola z reguły & Kod zgodny z tablicą \\
\bottomrule
\end{longtable}
}
Warunki końcowe: repozytorium usunięte; repo prywatne nie ujawniło swojego istnienia (404, nie 403).

\subsubsection{TC-F15 – Nadanie i odebranie dostępu kolaboratorowi}
{\small
\begin{longtable}{p{\dimexpr 0.270\linewidth-2\tabcolsep\relax}p{\dimexpr 0.730\linewidth-2\tabcolsep\relax}}
\toprule
\textbf{Pole} & \textbf{Wartość} \\
\midrule
\endhead
Priorytet & Wysoki \\
Technika & Testowanie przejść między stanami (brak dostępu $\rightarrow$ zapis $\rightarrow$ brak dostępu) \\
Typ testu & Funkcjonalny \\
Poziom testu & Testowanie systemowe (REST API) \\
Warunki wstępne / dane & Prywatne repo właściciela; drugi użytkownik \\
Powiązane funkcje & F07 \\
\bottomrule
\end{longtable}
}
{\small
\begin{longtable}{p{\dimexpr 0.050\linewidth-2\tabcolsep\relax}p{\dimexpr 0.480\linewidth-2\tabcolsep\relax}p{\dimexpr 0.470\linewidth-2\tabcolsep\relax}}
\toprule
\textbf{\#} & \textbf{Krok} & \textbf{Oczekiwany rezultat} \\
\midrule
\endhead
1 & Drugi użytkownik: \texttt{GET /repos/<repo>} & 404 \\
2 & Właściciel: \texttt{PUT /repos/<repo>/collaborators/<drugi>} z \texttt{permission: write} & 204 \\
3 & Drugi użytkownik: \texttt{GET /repos/<repo>} & 200 \\
4 & Drugi użytkownik: \texttt{POST /repos/<repo>/issues} & 201 – uprawnienie zapisu działa \\
5 & Właściciel: \texttt{DELETE .../collaborators/<drugi>}; drugi: \texttt{GET} & 404 \\
\bottomrule
\end{longtable}
}
Warunki końcowe: drugi użytkownik nie ma dostępu; repo usunięte.

\subsubsection{TC-F16 – Licznik postępu kamienia milowego}
{\small
\begin{longtable}{p{\dimexpr 0.270\linewidth-2\tabcolsep\relax}p{\dimexpr 0.730\linewidth-2\tabcolsep\relax}}
\toprule
\textbf{Pole} & \textbf{Wartość} \\
\midrule
\endhead
Priorytet & Średni \\
Technika & Testowanie oparte na kryteriach akceptacji \\
Typ testu & Funkcjonalny \\
Poziom testu & Testowanie systemowe (REST API) \\
Warunki wstępne / dane & Repo użytkownika; milestone „v1.0”; 2 zgłoszenia przypisane do milestone \\
Powiązane funkcje & F09, F05 \\
\bottomrule
\end{longtable}
}
{\small
\begin{longtable}{p{\dimexpr 0.050\linewidth-2\tabcolsep\relax}p{\dimexpr 0.480\linewidth-2\tabcolsep\relax}p{\dimexpr 0.470\linewidth-2\tabcolsep\relax}}
\toprule
\textbf{\#} & \textbf{Krok} & \textbf{Oczekiwany rezultat} \\
\midrule
\endhead
1 & Utwórz milestone i 2 zgłoszenia z \texttt{milestone} & 201 dla każdego \\
2 & \texttt{GET /milestones/<id>} & \texttt{open\_issues} = 2, \texttt{closed\_issues} = 0 \\
3 & Zamknij jedno zgłoszenie (\texttt{PATCH state=closed}) & 201 \\
4 & \texttt{GET /milestones/<id>} & \texttt{open\_issues} = 1, \texttt{closed\_issues} = 1 \\
\bottomrule
\end{longtable}
}
Warunki końcowe: milestone z postępem 50\%.

\subsubsection{TC-F17 – Uwierzytelnianie tokenem API}
{\small
\begin{longtable}{p{\dimexpr 0.270\linewidth-2\tabcolsep\relax}p{\dimexpr 0.730\linewidth-2\tabcolsep\relax}}
\toprule
\textbf{Pole} & \textbf{Wartość} \\
\midrule
\endhead
Priorytet & Średni \\
Technika & Podział na klasy równoważności (token poprawny / brak / błędny / unieważniony) \\
Typ testu & Funkcjonalny (bezpieczeństwo – uwierzytelnianie) \\
Poziom testu & Testowanie systemowe (REST API) \\
Warunki wstępne / dane & Użytkownik testowy; token \texttt{tok-<losowy>} z zakresem \texttt{read:user} \\
Powiązane funkcje & F12 \\
\bottomrule
\end{longtable}
}
{\small
\begin{longtable}{p{\dimexpr 0.050\linewidth-2\tabcolsep\relax}p{\dimexpr 0.480\linewidth-2\tabcolsep\relax}p{\dimexpr 0.470\linewidth-2\tabcolsep\relax}}
\toprule
\textbf{\#} & \textbf{Krok} & \textbf{Oczekiwany rezultat} \\
\midrule
\endhead
1 & \texttt{POST /users/<user>/tokens} & 201, zwrócony \texttt{sha1} \\
2 & \texttt{GET /user} z tokenem poprawnym & 200 \\
3 & \texttt{GET /user} bez nagłówka \texttt{Authorization} & 401 \\
4 & \texttt{GET /user} z tokenem \texttt{000\dots{}0} & 401 \\
5 & Usuń token, \texttt{GET /user} z usuniętym tokenem & 401 \\
\bottomrule
\end{longtable}
}
Warunki końcowe: token usunięty.

\subsection{Przypadki niefunkcjonalne (ISO/IEC 25010:2023)}
\subsubsection{TC-NF01 – Czas odpowiedzi pod obciążeniem}
{\small
\begin{longtable}{p{\dimexpr 0.270\linewidth-2\tabcolsep\relax}p{\dimexpr 0.730\linewidth-2\tabcolsep\relax}}
\toprule
\textbf{Pole} & \textbf{Wartość} \\
\midrule
\endhead
Priorytet & Wysoki \\
Charakterystyka & Wydajność (Performance efficiency) $\rightarrow$ zachowanie w czasie, wykorzystanie zasobów (przepustowość) \\
Technika & Testowanie obciążeniowe z profilem operacyjnym (mieszanka odczytów 3:3:2:2:1) \\
Typ testu & Niefunkcjonalny – wydajnościowy \\
Poziom testu & Testowanie systemowe \\
Warunki wstępne / dane & Repo \texttt{perf-target} z 20 zgłoszeniami; 20 wirtualnych użytkowników, narastanie 5/s, czas 30 s (CI: 15 s); narzędzie Locust 2.32 \\
Powiązane funkcje & F04, F05, F10 \\
\bottomrule
\end{longtable}
}
{\small
\begin{longtable}{p{\dimexpr 0.050\linewidth-2\tabcolsep\relax}p{\dimexpr 0.480\linewidth-2\tabcolsep\relax}p{\dimexpr 0.470\linewidth-2\tabcolsep\relax}}
\toprule
\textbf{\#} & \textbf{Krok} & \textbf{Oczekiwany rezultat} \\
\midrule
\endhead
1 & Przygotuj dane testowe przez API & Repo i 20 zgłoszeń utworzone \\
2 & Uruchom Locust w trybie headless z profilem \texttt{perf/locustfile.py} & Wygenerowany \texttt{reports/perf/locust\_stats.csv} \\
3 & Odczytaj statystyki zagregowane & Odsetek błędów = 0\% \\
4 & Odczytaj przepustowość & $\geq$ 10 żądań/s \\
5 & Odczytaj 95. percentyl każdego typu żądania & p95 < 500 ms dla każdego \\
\bottomrule
\end{longtable}
}
Warunki końcowe: dane wydajnościowe zapisane w \texttt{reports/perf/summary.json}; repo usunięte.

\subsubsection{TC-NF02 – Czas ładowania kluczowych stron}
{\small
\begin{longtable}{p{\dimexpr 0.270\linewidth-2\tabcolsep\relax}p{\dimexpr 0.730\linewidth-2\tabcolsep\relax}}
\toprule
\textbf{Pole} & \textbf{Wartość} \\
\midrule
\endhead
Priorytet & Średni \\
Charakterystyka & Wydajność $\rightarrow$ zachowanie w czasie \\
Technika & Pomiar z użyciem Navigation Timing API (mediana z 3 próbek) \\
Typ testu & Niefunkcjonalny – wydajnościowy \\
Poziom testu & Testowanie systemowe (Selenium) \\
Warunki wstępne / dane & Zalogowany użytkownik, repo testowe; strony: \texttt{/}, \texttt{/explore/repos}, repo, lista zgłoszeń, \texttt{/user/settings} \\
Powiązane funkcje & F02, F04, F05, F10 \\
\bottomrule
\end{longtable}
}
{\small
\begin{longtable}{p{\dimexpr 0.050\linewidth-2\tabcolsep\relax}p{\dimexpr 0.480\linewidth-2\tabcolsep\relax}p{\dimexpr 0.470\linewidth-2\tabcolsep\relax}}
\toprule
\textbf{\#} & \textbf{Krok} & \textbf{Oczekiwany rezultat} \\
\midrule
\endhead
1 & Zaloguj się w przeglądarce & Sesja aktywna \\
2 & Dla każdej strony 3$\times$ otwórz i odczytaj \texttt{loadEventEnd - startTime} & Pomiar dostępny \\
3 & Policz medianę dla każdej strony & Mediana < 2000 ms dla każdej strony \\
\bottomrule
\end{longtable}
}
Warunki końcowe: wyniki w \texttt{reports/page\_load.json}.

\subsubsection{TC-NF03 – Zgodność kluczowych widoków z WCAG 2.1 AA}
{\small
\begin{longtable}{p{\dimexpr 0.270\linewidth-2\tabcolsep\relax}p{\dimexpr 0.730\linewidth-2\tabcolsep\relax}}
\toprule
\textbf{Pole} & \textbf{Wartość} \\
\midrule
\endhead
Priorytet & Średni \\
Charakterystyka & Zdolność do interakcji (Interaction capability) $\rightarrow$ inkluzywność / dostępność \\
Technika & Testowanie oparte na liście kontrolnej (reguły axe-core 4.10.2: wcag2a, wcag2aa, wcag21aa) \\
Typ testu & Niefunkcjonalny – dostępności \\
Poziom testu & Testowanie systemowe (Selenium + axe-core) \\
Warunki wstępne / dane & Widoki: logowanie, rejestracja, repo, lista zgłoszeń, nowe repo \\
Powiązane funkcje & F01, F02, F04, F05 \\
\bottomrule
\end{longtable}
}
{\small
\begin{longtable}{p{\dimexpr 0.050\linewidth-2\tabcolsep\relax}p{\dimexpr 0.480\linewidth-2\tabcolsep\relax}p{\dimexpr 0.470\linewidth-2\tabcolsep\relax}}
\toprule
\textbf{\#} & \textbf{Krok} & \textbf{Oczekiwany rezultat} \\
\midrule
\endhead
1 & Otwórz widoki anonimowe i uruchom axe & Raport naruszeń dla każdego widoku \\
2 & Zaloguj się, otwórz widoki wymagające sesji i uruchom axe & Raport naruszeń dla każdego widoku \\
3 & Zapisz podsumowanie do \texttt{reports/axe.json} & Plik zawiera id, wpływ i liczbę węzłów \\
4 & Sprawdź naruszenia o wpływie \texttt{critical} & Brak naruszeń krytycznych \\
\bottomrule
\end{longtable}
}
Warunki końcowe: raport dostępności zapisany. \textbf{Wynik: niezaliczony – DEF-02.}

\subsubsection{TC-NF04 – Podstawowe zabezpieczenia aplikacji webowej}
{\small
\begin{longtable}{p{\dimexpr 0.270\linewidth-2\tabcolsep\relax}p{\dimexpr 0.730\linewidth-2\tabcolsep\relax}}
\toprule
\textbf{Pole} & \textbf{Wartość} \\
\midrule
\endhead
Priorytet & Wysoki \\
Charakterystyka & Bezpieczeństwo (Security) $\rightarrow$ poufność, integralność \\
Technika & Testowanie oparte na liście kontrolnej (OWASP ASVS: V3 sesje, V4 kontrola dostępu, V14 nagłówki) \\
Typ testu & Niefunkcjonalny – bezpieczeństwa \\
Poziom testu & Testowanie systemowe (HTTP) \\
Warunki wstępne / dane & Użytkownik testowy; obca domena \texttt{http://evil.example} \\
Powiązane funkcje & F02, F07 \\
\bottomrule
\end{longtable}
}
{\small
\begin{longtable}{p{\dimexpr 0.050\linewidth-2\tabcolsep\relax}p{\dimexpr 0.480\linewidth-2\tabcolsep\relax}p{\dimexpr 0.470\linewidth-2\tabcolsep\relax}}
\toprule
\textbf{\#} & \textbf{Krok} & \textbf{Oczekiwany rezultat} \\
\midrule
\endhead
1 & \texttt{GET /user/login}, sprawdź nagłówki & \texttt{X-Content-Type-Options: nosniff}, \texttt{X-Frame-Options: SAMEORIGIN} lub \texttt{DENY} \\
2 & Zaloguj się, sprawdź \texttt{Set-Cookie} ciasteczka sesji & \texttt{HttpOnly} i \texttt{SameSite=Lax} \\
3 & Jako zalogowany wyślij \texttt{POST /repo/create} z \texttt{Origin: evil.example}, \texttt{Sec-Fetch-Site: cross-site} & 403; repo nie powstało \\
4 & Wyślij \texttt{POST /user/login} cross-site & 403 (\textbf{faktycznie 303 – DEF-01, oznaczony jako xfail}) \\
5 & Anonimowo \texttt{GET} strony i surowego pliku repo prywatnego & 404 dla obu \\
\bottomrule
\end{longtable}
}
Warunki końcowe: brak utworzonych zasobów.

\subsubsection{TC-NF05 – Kompatybilność przeglądarek}
{\small
\begin{longtable}{p{\dimexpr 0.270\linewidth-2\tabcolsep\relax}p{\dimexpr 0.730\linewidth-2\tabcolsep\relax}}
\toprule
\textbf{Pole} & \textbf{Wartość} \\
\midrule
\endhead
Priorytet & Średni \\
Charakterystyka & Kompatybilność (Compatibility) $\rightarrow$ współistnienie / interoperacyjność; Elastyczność (Flexibility) $\rightarrow$ adaptowalność \\
Technika & Testowanie sparametryzowane (ten sam scenariusz w wielu konfiguracjach) \\
Typ testu & Niefunkcjonalny – kompatybilności \\
Poziom testu & Testowanie systemowe (Selenium) \\
Warunki wstępne / dane & Google Chrome 157, Microsoft Edge (lokalnie); w CI tylko Chrome \\
Powiązane funkcje & F02, F03 \\
\bottomrule
\end{longtable}
}
{\small
\begin{longtable}{p{\dimexpr 0.050\linewidth-2\tabcolsep\relax}p{\dimexpr 0.480\linewidth-2\tabcolsep\relax}p{\dimexpr 0.470\linewidth-2\tabcolsep\relax}}
\toprule
\textbf{\#} & \textbf{Krok} & \textbf{Oczekiwany rezultat} \\
\midrule
\endhead
1 & Uruchom przeglądarkę i zaloguj się & Awatar w nawigacji \\
2 & Utwórz repozytorium z README & Lista plików zawiera \texttt{README.md} \\
3 & Wyloguj się & Link „Sign In” w nawigacji \\
\bottomrule
\end{longtable}
}
Warunki końcowe: identyczny wynik w każdej przeglądarce; repo usunięte.

